% ─────────────────────────────────────────────────────────────────────────────
\section{Periodic Simple Continued Fractions}
% ─────────────────────────────────────────────────────────────────────────────

By Lagrange's theorem, a real number has an eventually periodic continued
fraction expansion if and only if it is a \emph{quadratic irrational} — a root
of a degree-two polynomial with integer coefficients that is not itself
rational.  The favorite examples are $\sqrt{2} = [1;\,\overline{2}]$,
$\varphi = [1;\,\overline{1}]$, and more generally every expression of the
form $(P + \sqrt{D})/Q$ with $D$ not a perfect square.
This section introduces \texttt{PeriodicSimpleContinuedFraction}, the \catena
class that represents such numbers exactly and exposes their structure through
the same functional interface used by \texttt{SimpleContinuedFraction}.
The deeper theory — the proof of Lagrange's theorem, the correspondence with
Pell equations, and the algebraic operations on quadratic irrationals — is
deferred to a forthcoming section.

\subsection{Purely periodic vs.\ eventually periodic}

\begin{definition}[Eventually periodic SCF]
  An infinite SCF $(a_0, g)$ with $g: \N \to \Zp$ is
  \emph{eventually periodic} if there exist integers $m \geq 0$ and $r \geq 1$
  such that
  \[
    g(m + k) = g(m + k + r) \qquad \text{for all } k \geq 0.
  \]
  The smallest such $r$ is the \emph{minimal period length} and the
  corresponding block $(g(m), \ldots, g(m+r-1))$ is the \emph{period}.
  The block $(g(0), \ldots, g(m-1))$ (empty when $m = 0$) is the
  \emph{pre-period}.
  When $m = 0$ the SCF is \emph{purely periodic} (its body repeats from the
  very first partial quotient).
\end{definition}

In \catena there is no distinction between the two at the type level.
Both are represented by a single class, \texttt{PeriodicSimpleContinuedFraction},
instantiated with the same constructor.  The difference is entirely
encoded by the pre-period length $m$: a purely periodic PSCF is simply
one where $m = 0$, so that $g_{\mathrm{pre}}$ has an empty domain and
contributes no partial quotients.  In Definition~\ref{def:pscf} this means
the case-split collapses to $g(n) = g_{\mathrm{per}}(n \bmod r)$ for all
$n \geq 0$, and the \texttt{pre\_period} argument to the constructor is
omitted (or passed as an empty sequence).

\begin{remark}
  Some sources define ``purely periodic'' to include $a_0$ in the repeating
  block, writing $[\overline{a_0;\,a_1,\ldots,a_{r-1}}]$.  In \catena the
  integer part $a_0$ is always stored separately and is never part of a
  periodic block.  ``Purely periodic'' here refers to the body $g$ alone:
  $m = 0$ means the body repeats from index~$0$, regardless of the value
  of $a_0$.
\end{remark}

\subsection{The integer-part separation and its effect on the period}

Because $a_0$ is stored outside the generator, the period reported by
\texttt{PeriodicSimpleContinuedFraction} is always the period of the
\emph{body generator} $g$, not of the full sequence $(a_0, a_1, a_2, \ldots)$.
This can shift the apparent repeating block relative to a naive reading of the
expansion.

\begin{example}
  Consider $x = [3;\,\overline{1,3}]$.  Written out, the partial quotients are
  $3,\,1,\,3,\,1,\,3,\ldots$  The digits $3$ and $1$ alternate, so one might
  expect the period to be $(3, 1)$.  However, $a_0 = 3$ is the integer part and
  is stored separately; the body sequence is $1, 3, 1, 3, \ldots$, which has
  period $(1, 3)$.  The number $[\overline{3,1}]$ (in the sense of a purely
  periodic expansion starting from $a_0$) is the same value — the two
  representations differ only in where the period is ``cut'' — but \catena
  always extracts $a_0 = \lfloor x \rfloor$ first, so the stored period is
  $(1, 3)$ in both cases.
\end{example}

This is not a limitation but a deliberate consequence of the uniform
representation: the integer part is always an integer, the body generator
always maps into $\Zp$, and the two are never mixed.

\subsection{Definition and data model}

\begin{definition}[Periodic simple continued fraction]
  \label{def:pscf}
  A \emph{periodic simple continued fraction} (PSCF) is an eventually
  periodic SCF $(a_0, g)$ whose body generator $g: \N \to \Zp$ is determined
  by two finite generators
  \begin{align*}
    g_{\mathrm{pre}} &: m \longrightarrow \Zp, \\
    g_{\mathrm{per}} &: r \longrightarrow \Zp, \quad r \geq 1,
  \end{align*}
  via the rule
  \[
    g(n) \;=\;
    \begin{cases}
      g_{\mathrm{pre}}(n)                              & \text{if } n < m, \\[4pt]
      g_{\mathrm{per}}\!\bigl((n - m) \bmod r\bigr)   & \text{if } n \geq m.
    \end{cases}
  \]
  Here $m \in \N$ is the \emph{pre-period length} and $r \in \Zp$ is the
  \emph{period length}.
  The sequences
  \[
    \bigl(g_{\mathrm{pre}}(0),\ldots,g_{\mathrm{pre}}(m-1)\bigr)
    \qquad\text{and}\qquad
    \bigl(g_{\mathrm{per}}(0),\ldots,g_{\mathrm{per}}(r-1)\bigr)
  \]
  are referred to as the \emph{pre-period} and the \emph{period} of the PSCF
  respectively.
\end{definition}

In \catena, $g_{\mathrm{pre}}$ and $g_{\mathrm{per}}$ are each stored as a
\texttt{FiniteGenerator} — a fixed, immutable array of positive integers backed
by a compact \texttt{array.array}.  These two \texttt{FiniteGenerator} instances
are then combined into a single \texttt{PeriodicGenerator}, which implements
the case-split of Definition~\ref{def:pscf} as its underlying callable.
\texttt{PeriodicSimpleContinuedFraction} wraps this generator together with the
integer part $a_0$, and exposes the two sub-generators through the
\texttt{non\_repeating\_part} and \texttt{period} properties.

\begin{catenadef}[\texttt{PeriodicSimpleContinuedFraction}]
\begin{lstlisting}[style=python]
from catena import PeriodicSimpleContinuedFraction as PSCF

# phi = [1; 1,1,1,...] -- golden ratio, purely periodic body (m = 0)
phi = PSCF(period=[1], integer_part=1)
print(phi.integer_part)        # 1
print(phi.period)              # (1,)
print(phi.non_repeating_part)  # ()

# sqrt(2) = [1; 2,2,2,...] -- purely periodic body (m = 0)
sqrt2 = PSCF(period=[2], integer_part=1)

# [3; 1,3,1,3,...] -- the body period is (1,3), not (3,1)
x = PSCF(period=[1, 3], integer_part=3)
print(x.period)              # (1, 3)
print(x.non_repeating_part)  # ()

# (7 + sqrt(10))/4 = [2; 1,1, 5,1,1,1,24,1,1,1,5,...]
# non-repeating pre-period (1,1) followed by period (5,1,1,1,24,1,1,1)
y = PSCF(period=[5, 1, 1, 1, 24, 1, 1, 1], pre_period=[1, 1], integer_part=2)
print(y.non_repeating_part)  # (1, 1)
print(y.period)              # (5, 1, 1, 1, 24, 1, 1, 1)
\end{lstlisting}
\end{catenadef}

\subsection{Inherited and additional methods}

\texttt{PeriodicSimpleContinuedFraction} inherits all of
\texttt{SimpleContinuedFraction} without change.  In particular,
\texttt{convergent}, \texttt{tail\_convergent}, \texttt{tail},
\texttt{segment}, \texttt{shift}, and \texttt{\_\_add\_\_} all work on a
PSCF exactly as described in the SCF section, since the underlying
\texttt{PeriodicGenerator} is just another callable $g: \N \to \Zp$.
The convergent cache and the tail-sharing machinery are likewise inherited
unchanged.

The following members are added or overridden by the class.

\begin{description}[leftmargin=2em, style=nextline]
  \item[\texttt{non\_repeating\_part}, \texttt{period}]
    Read-only properties returning the pre-period
    $(g_{\mathrm{pre}}(0),\ldots,g_{\mathrm{pre}}(m-1))$ and the period
    $(g_{\mathrm{per}}(0),\ldots,g_{\mathrm{per}}(r-1))$ as immutable
    tuples.

  \item[\texttt{from\_quadratic\_surd(P, Q, D)}]
    Class method.  Constructs a PSCF whose value is $(P + \sqrt{D})/Q$ by
    running the quadratic-surd expansion algorithm and forwarding the
    resulting pre-period, period, and integer part to the constructor.
    This is the factory used internally by \texttt{inverse} and
    \texttt{conjugate}.

  \item[\texttt{\_\_float\_\_} (override)]
    Returns \texttt{(P + D**0.5) / Q} from the cached
    \texttt{quadratic\_surd()} triple.  Unlike
    \texttt{SimpleContinuedFraction.\_\_float\_\_}, which evaluates
    \texttt{convergent(50)}, the periodic override is exact up to
    IEEE\,754 rounding in a single arithmetic expression, exploiting the
    closed-form surd representation.

  \item[\texttt{as\_decimal()}]
    Returns the value as an arbitrary-precision \texttt{Decimal} via
    \texttt{Decimal(D).sqrt()}.  This is the periodic analogue of
    \texttt{FiniteSimpleContinuedFraction.to\_decimal()}: both are exact
    up to the current \texttt{decimal} context precision, but here
    exactness comes from the algebraic form rather than from evaluating the
    terminal convergent.
\end{description}

The following new methods are also implemented but their correctness rests on
algebraic properties of quadratic irrationals that are not developed in
this version of the document.  They are listed here for completeness;
full justification is deferred to a forthcoming section.

\begin{description}[leftmargin=2em, style=nextline]
  \item[\texttt{quadratic\_coefficients()}]
    Returns $(A, B, C)$ with $\gcd(A,B,C) = 1$ and $A > 0$ such that the
    value of the PSCF satisfies $Ax^2 + Bx + C = 0$.
  \item[\texttt{quadratic\_surd()}]
    Returns $(P, Q, D)$ such that the value equals $(P + \sqrt{D})/Q$,
    derived from the quadratic coefficients.
  \item[\texttt{is\_principal\_surd()}, \texttt{is\_conjugate\_root()}]
    Predicates distinguishing the two roots of the minimal polynomial,
    identified by the sign of $Q$.
  \item[\texttt{inverse()} (override)]
    Returns $1/x$ as a new PSCF, computed via the surd formula.  Overrides
    the base-class implementation to exploit the closed form; the
    write-once involution contract (\texttt{\_inverse.\_inverse is self})
    is preserved.
  \item[\texttt{conjugate()}]
    Returns the algebraic conjugate $(P - \sqrt{D})/Q$, i.e.\ the other
    root of the same minimal polynomial.
  \item[\texttt{\_\_neg\_\_()}]
    Returns $-x$ as a new PSCF via the surd negation $(P,\,-Q,\,D)$. The closed-form surd 
    representation ensures that the definition of negation is exact, so no approximation is involved.
\end{description}
